% ============================================================
% CAPITOLO 6 - IL COLLEGAMENTO ARMONICO
% ============================================================
\chapter{Il collegamento armonico}
\label{chap:collegamento}

Collegare due accordi significa portare le quattro voci dal primo al secondo con il minimo sforzo e senza spezzare la linea di nessuna parte.

\section{Collegare due accordi}
\label{sec:collegamento_armonico}

Il principio guida è la \textbf{legge del minimo sforzo}: le voci si muovono il meno possibile per raggiungere l'accordo successivo.

\begin{regola}[Principio di minima azione]
	\label{reg:minima_azione}
	Per collegare due accordi si seguono due criteri, in questo ordine:
	\begin{enumerate}
		\item \textbf{Note comuni:} se i due accordi hanno una o più note in comune, queste restano ferme nella stessa voce (moto obliquo).
		\item \textbf{Grado congiunto:} le altre voci seguono il percorso più breve, preferibilmente per grado congiunto (intervalli di seconda).
	\end{enumerate}
\end{regola}

\section{Collegamenti per distanza delle fondamentali}
\label{sec:collegamento_distanza}

Il numero di note comuni dipende dalla distanza tra le fondamentali dei due accordi.

\begin{itemize}
	\item \textbf{Relazione di quinta (o quarta):} gli accordi hanno \textbf{una nota in comune}. Esempi: I-IV, I-V, V-I.
	\item \textbf{Relazione di terza (o sesta):} gli accordi hanno \textbf{due note in comune}. È il collegamento più fluido. Esempi: I-VI, I-III.
	\item \textbf{Relazione di seconda (o settima):} gli accordi \textbf{non hanno note in comune}. Per evitare parallelismi le voci superiori procedono per moto contrario rispetto al basso. Esempi: IV-V, V-VI.
\end{itemize}

\section{La risoluzione della sensibile}
\label{sec:risoluzione_sensibile}

La \textbf{sensibile}, VII grado della scala, è il grado più direzionato del sistema tonale: dista un semitono dalla tonica e la sua attrazione va rispettata nella condotta delle voci.

\begin{regola}[I tre movimenti possibili della sensibile]
	\label{reg:tre_movimenti_sensibile}
	La sensibile può muoversi in tre modi:
	\begin{enumerate}
		\item \textbf{Sale di semitono alla tonica.} È la risoluzione primaria, obbligatoria se la sensibile è al soprano.
		\item \textbf{Scende di terza alla dominante}, a condizione che si trovi in una \textbf{parte interna} (contralto o tenore) e che il basso proceda in \textbf{moto contrario} rispetto a essa. Così si ottiene una tonica completa, senza triplicare la fondamentale.
		\item \textbf{Scende di semitono cromatico}, solo in presenza di un \textbf{cambio di tonalità o di una tonicizzazione}: la sensibile abbassata non è più sensibile della tonalità d'impianto ma un altro grado della tonalità locale.
	\end{enumerate}
	Ogni altro movimento va evitato, salvo che la sensibile sia in una voce intermedia e il contesto lo giustifichi.
\end{regola}

\begin{esempio}[Le risoluzioni della sensibile in Do maggiore]
	\label{ex:risoluzioni_sensibile}

	\textbf{A. La sensibile sale alla tonica:}
	\musica{es_cap04_01}

	\vspace{0.5cm}
	\textbf{B. La sensibile scende di terza (voce interna, moto contrario con il basso; qui nel collegamento \sigla{V}{7}{}-I, studiato nel Capitolo \ref{chap:settima_dominante}):}
	\musica{es_cap04_02}

	Nel caso B il Si sta al contralto e scende di terza al Sol, mentre il basso sale da Sol a Do. Il moto contrario e la posizione interna rendono corretto il collegamento, e la tonica risulta completa.
\end{esempio}

\begin{errore}[Sensibile non risolta nel soprano]
	\label{err:sensibile_non_risolta}
	Se la sensibile è al \textbf{soprano} deve salire di semitono alla tonica. Il soprano non può scendere di terza, anche se la nota di arrivo appartiene all'accordo: la linea risulterebbe poco cantabile e armonicamente incoerente.
\end{errore}
